\begin{abstract}
\noindent
AWS exposes the same workload to three very different compute paradigms ---
virtual machines on EC2, managed containers on ECS Fargate, and serverless
functions on Lambda. Selecting between them is rarely a clean choice because
performance, elasticity, and cost rarely move in the same direction. We
deploy an identical container image of a URL-shortener service backed by
DynamoDB to all three paradigms behind a single Application Load Balancer,
and benchmark them against four scenarios: a steady-state RPS sweep, a
step-arrival burst, a cold-start probe, and a mixed read/write workload.
To keep the comparison fair we use the AWS Lambda Web Adapter so that the
exact same FastAPI binary runs on every paradigm, and we drive load from a
co-located k6 worker running in the same availability zone. Over the
regime we could sustain within a new AWS account's default service
quotas, Lambda delivered the best and most stable latency
(p95\,$\approx$\,40\,ms and p99\,$\approx$\,52\,ms across 25--75\,RPS),
a single-instance EC2~ASG held a 13\,ms median up to 200\,RPS before its
burstable-credit tail blew out, and a single 0.5\,vCPU Fargate task
saturated by 100\,RPS with p50 above 3\,s. We combine these measurements
with public April~2026 us-east-1 pricing to derive a break-even crossover
at roughly \textbf{20\,RPS}: below that rate Lambda's no-idle billing is
cheapest; above it, pre-provisioned compute wins. Our infrastructure is
fully described in Terraform, the k6 scripts and analysis pipeline
reproduce every figure from raw data with a single \texttt{make} target,
and the artifact repository records the end-to-end process including the
migration from the AWS Academy Learner Lab (whose sandboxed account was
deactivated mid-experiment) to a personal AWS account.
\end{abstract}
